\documentclass{beamer}
\usepackage{amssymb}
\usepackage{tikz}

\begin{document}
  \begin{frame}
    \textit{Temas}:

    \begin{itemize}
      \item Interés compuesto
      \item Valor futuro y valor presente
      \item Tasas de interés y tasas equivalentes
      \item Flujo de caja
      \item Equivalencia financiera y ecuación de valor
      \item Tasa de interés variable
      \item Sustitución y refinanciación de obligaciones
    \end{itemize}
  \end{frame}

  \begin{frame}
    \textit{Definición}: En el interés compuesto, los intereses generados en cada
    periodo se suman al capital y producen, a su vez, intereses en los periodos
    siguientes. A este proceso se le llama capitalización.
  \end{frame}

  \begin{frame}
    \textit{Definición}: En el interés compuesto, los intereses generados en cada
    periodo se suman al capital y producen, a su vez, intereses en los periodos
    siguientes. A este proceso se le llama capitalización.
    \newline\newline

    Sean $P$ el capital inicial, $i$ la tasa de interés por periodo y $n$ el
    número de periodos.
  \end{frame}

  \begin{frame}
    Sean $P$ el capital inicial, $i$ la tasa de interés por periodo y $n$ el
    número de periodos.

    \begin{equation*}
      \begin{aligned}
        F_1 &= P + Pi = P(1+i) \\
        F_2 &= P(1+i) + P(1+i)i = P(1+i)^2 \\
        &\;\;\vdots \\
        F_n &= P(1+i)^n
      \end{aligned}
    \end{equation*}
  \end{frame}

  \begin{frame}
    \textit{Fórmula}: El valor futuro de un capital $P$ después de $n$ periodos
    a la tasa $i$ es
    \[F = P(1+i)^n\]
  \end{frame}

  \begin{frame}
    \textit{Fórmula}: El valor futuro de un capital $P$ después de $n$ periodos
    a la tasa $i$ es
    \[F = P(1+i)^n\]

    \textit{Fórmula}: El valor presente de una cantidad $F$ que se recibe dentro
    de $n$ periodos es
    \[P = F(1+i)^{-n}\]
  \end{frame}

  \begin{frame}
    \textit{Fórmula}: El valor futuro de un capital $P$ después de $n$ periodos
    a la tasa $i$ es
    \[F = P(1+i)^n\]

    \textit{Fórmula}: El valor presente de una cantidad $F$ que se recibe dentro
    de $n$ periodos es
    \[P = F(1+i)^{-n}\]

    El dinero tiene valor en el tiempo: llevar un valor al futuro es
    capitalizar y llevarlo al presente es descontar.
  \end{frame}

  \begin{frame}
    \textit{Definición}: Una tasa del $j\%$ m.v. (mensual vencida) es una tasa
    que se aplica cada mes y cuyos intereses se causan al final de cada mes.
    \newline\newline

    El periodo de la tasa debe coincidir con la unidad en que se mide el
    tiempo: si la tasa es mensual, $n$ se cuenta en meses.
  \end{frame}

  \begin{frame}
    \textit{Definición}: Una tasa del $j\%$ m.v. (mensual vencida) es una tasa
    que se aplica cada mes y cuyos intereses se causan al final de cada mes.
    \newline\newline

    El periodo de la tasa debe coincidir con la unidad en que se mide el
    tiempo: si la tasa es mensual, $n$ se cuenta en meses.
    \newline\newline

    Si $i$ es la tasa mensual, la tasa nominal anual capitalizable mensualmente
    es $12i$.
  \end{frame}

  \begin{frame}
    \textit{Definición}: Dos tasas con periodos distintos son equivalentes si
    producen el mismo valor futuro sobre un mismo capital y en el mismo tiempo.
    \newline\newline

    Si $i_1$ se capitaliza $m_1$ veces y $i_2$ se capitaliza $m_2$ veces al año,
    \[(1+i_1)^{m_1} = (1+i_2)^{m_2}\]
  \end{frame}

  \begin{frame}
    \textit{Definición}: Dos tasas con periodos distintos son equivalentes si
    producen el mismo valor futuro sobre un mismo capital y en el mismo tiempo.
    \newline\newline

    Si $i_1$ se capitaliza $m_1$ veces y $i_2$ se capitaliza $m_2$ veces al año,
    \[(1+i_1)^{m_1} = (1+i_2)^{m_2}\]

    En particular, la tasa efectiva anual de una tasa mensual $i$ es
    \[i_{ea} = (1+i)^{12} - 1\]
  \end{frame}

  \begin{frame}
    \textit{Definición}: Un diagrama de flujo de caja (línea de tiempo) es la
    representación, sobre un eje de tiempo, de los ingresos y egresos de una
    operación financiera.
  \end{frame}

  \begin{frame}
    \textit{Definición}: Un diagrama de flujo de caja (línea de tiempo) es la
    representación, sobre un eje de tiempo, de los ingresos y egresos de una
    operación financiera.

    \begin{center}
      \begin{tikzpicture}[scale=0.9]
        \draw[->] (0,0) -- (6.4,0) node[right] {\small $t$};
        \foreach \x/\l in {0/0, 2/$n_1$, 4/$n_2$, 6/$n_3$}
          \draw (\x,0.08) -- (\x,-0.08) node[below] {\small \l};
        \draw[->, thick] (0,0.1) -- (0,1);
        \draw[->, thick] (2,-0.7) -- (2,-1.5);
        \draw[->, thick] (4,-0.7) -- (4,-1.5);
        \draw[->, thick] (6,-0.7) -- (6,-1.5);
      \end{tikzpicture}
    \end{center}
  \end{frame}

  \begin{frame}
    \textit{Definición}: Un diagrama de flujo de caja (línea de tiempo) es la
    representación, sobre un eje de tiempo, de los ingresos y egresos de una
    operación financiera.

    \begin{itemize}
      \item El momento $t = 0$ corresponde a ``hoy''.
      \item Las flechas indican el sentido del flujo (ingreso o egreso).
      \item Cada pago se ubica en el periodo en que ocurre.
    \end{itemize}
  \end{frame}

  \begin{frame}
    \textit{Definición}: Dos conjuntos de pagos son equivalentes a la tasa $i$ si,
    llevados a una misma fecha, tienen el mismo valor. Esa fecha se llama fecha
    focal.
  \end{frame}

  \begin{frame}
    \textit{Definición}: Dos conjuntos de pagos son equivalentes a la tasa $i$ si,
    llevados a una misma fecha, tienen el mismo valor. Esa fecha se llama fecha
    focal.
    \newline\newline

    Un pago $C$ que ocurre en el periodo $t$ tiene, en la fecha focal $f$, el
    valor
    \[C(1+i)^{f-t}\]
  \end{frame}

  \begin{frame}
    \textit{Definición}: Dos conjuntos de pagos son equivalentes a la tasa $i$ si,
    llevados a una misma fecha, tienen el mismo valor. Esa fecha se llama fecha
    focal.
    \newline\newline

    Un pago $C$ que ocurre en el periodo $t$ tiene, en la fecha focal $f$, el
    valor
    \[C(1+i)^{f-t}\]

    \textit{Ecuación de valor}: En la fecha focal, la suma de los valores de las
    obligaciones es igual a la suma de los valores de los pagos.
  \end{frame}

  \begin{frame}
    \textit{Teorema}: Con una tasa de interés compuesto constante, la solución de
    una ecuación de valor no depende de la fecha focal elegida.
    \newline\newline

    Se escoge la fecha que simplifique las potencias, por ejemplo la fecha del
    pago desconocido.
  \end{frame}

  \begin{frame}
    \textit{Definición}: Cuando la tasa de interés cambia durante la operación, se
    aplica cada tasa en su respectivo tramo de tiempo.
  \end{frame}

  \begin{frame}
    \textit{Definición}: Cuando la tasa de interés cambia durante la operación, se
    aplica cada tasa en su respectivo tramo de tiempo.
    \newline\newline

    Si la tasa es $i_1$ durante $n_1$ periodos e $i_2$ durante $n_2$ periodos,
    \[F = P(1+i_1)^{n_1}(1+i_2)^{n_2}\]
  \end{frame}

  \begin{frame}
    \textit{Definición}: Cuando la tasa de interés cambia durante la operación, se
    aplica cada tasa en su respectivo tramo de tiempo.
    \newline\newline

    Si la tasa es $i_1$ durante $n_1$ periodos e $i_2$ durante $n_2$ periodos,
    \[F = P(1+i_1)^{n_1}(1+i_2)^{n_2}\]

    Para el valor presente, se descuenta tramo a tramo:
    \[P = F(1+i_2)^{-n_2}(1+i_1)^{-n_1}\]
  \end{frame}

  \begin{frame}
    \textit{Definición}: Sustituir (o refinanciar) una obligación es reemplazar un
    conjunto de pagos por otro equivalente, con nuevos montos y nuevas fechas.
  \end{frame}

  \begin{frame}
    \textit{Definición}: Sustituir (o refinanciar) una obligación es reemplazar un
    conjunto de pagos por otro equivalente, con nuevos montos y nuevas fechas.

    \begin{itemize}
      \item La obligación original se valora con su tasa.
    \end{itemize}
  \end{frame}

  \begin{frame}
    \textit{Definición}: Sustituir (o refinanciar) una obligación es reemplazar un
    conjunto de pagos por otro equivalente, con nuevos montos y nuevas fechas.

    \begin{itemize}
      \item La obligación original se valora con su tasa.
      \item Los nuevos pagos se valoran con la tasa del nuevo acuerdo, que puede
      ser distinta.
      \item Ambos valores se igualan en una misma fecha focal.
    \end{itemize}
  \end{frame}

  \begin{frame}
    \textit{Planteamiento}: Cuando un pago es desconocido, se expresa en función
    de una incógnita y se plantea la ecuación de valor. Algunas relaciones
    frecuentes son:

    \begin{itemize}
      \item Un pago es una fracción del valor de contado: $C = pV$.
    \end{itemize}
  \end{frame}

  \begin{frame}
    \textit{Planteamiento}: Cuando un pago es desconocido, se expresa en función
    de una incógnita y se plantea la ecuación de valor. Algunas relaciones
    frecuentes son:

    \begin{itemize}
      \item Un pago es una fracción del valor de contado: $C = pV$.
      \item Pagos iguales: $C_1 = C_2 = C$.
      \item Pagos en progresión geométrica de razón $r$: $C_k = r\,C_{k-1}$.
      \item Un saldo final igual al depósito inicial.
    \end{itemize}
  \end{frame}

  \begin{frame}
    \textit{Interpretación}: El resultado de una ecuación de valor es un valor
    en la fecha focal elegida y debe leerse en el contexto del problema.

    \begin{itemize}
      \item Verificar que el periodo de la tasa y el del tiempo coincidan.
    \end{itemize}
  \end{frame}

  \begin{frame}
    \textit{Interpretación}: El resultado de una ecuación de valor es un valor
    en la fecha focal elegida y debe leerse en el contexto del problema.

    \begin{itemize}
      \item Verificar que el periodo de la tasa y el del tiempo coincidan.
      \item Un valor presente debe ser menor que el valor futuro
      correspondiente.
      \item Indicar qué representa el resultado: cuánto, en qué fecha y a qué
      tasa.
    \end{itemize}
  \end{frame}
\end{document}
